\begin{homeworkProblem}[Seção 1-3]
  Sejam $V \subset U \subset \mathbb{R}^m$ abertos e $\delta > 0$ tais que $x \in V$, $|h|<\delta \Rightarrow x+h \in U$. Seja $B = B(0;\delta)$. Se $f : U \to \mathbb{R}^n$ é diferenciável, então $\varphi : V \times B \to \mathbb{R}^n$, definida por $\varphi(x,h)=f(x+h)$, é diferenciável, sendo
  \begin{align*}
    \varphi'(x_0,h_0)\cdot (u,v) = f'(x_0+h_0)\cdot (u+v).   
  \end{align*}
  \begin{solution}
    Como $f$ é diferenciável, sabemos que dado $x\in U$ y $v\in \mathbb{R}^{m}$ o suficientemente pequeno para que $x+v\in U$, então
    \begin{align*}
      f(x+v)=f(x)+f^{\prime}(x)\cdot v+r(v).
    \end{align*}
    Agora, seja $x,h\in V$ e $u,v\in B$, logo usando a definição de $\varphi$ e a difererenciabilidade do $f$ temos 
    \begin{align*}
      \varphi(x+u,h+v)&=f(x+u+h+v)\\
      &=f(x+h+u+v)\\
      &=f(x+h)+f^{\prime}(x+h)\cdot(u+v)+r(u+v)\\
      &=\varphi(x,y)+f^{\prime}(x+h)\cdot (u+v)+r(u+v),\quad\text{onde }\lim_{u+v \to 0}\frac{r(u+v)}{|u+v|}=0.
    \end{align*}
    Portanto, nos podemos afirmar que $\varphi^{\prime}(x,h)\cdot(u,v)=f^{\prime}(x+h)\cdot(u+v)$.
  \end{solution}
\end{homeworkProblem}
